\section{附录：社区平台变更评审表}

本附录把本讲的案例压成一套可复用评审流程，适用于内容政策、API、付费功能、AI 摘要和 agent 接口等高影响变更。评审不是为了让所有群体达成一致，而是要求团队在上线前公开关键依赖、成本、权力变化和失败恢复路径。每一项都应有负责人、证据和可停止条件，避免把“社区会适应”当作默认假设。

\begin{table}[H]
\centering
\begin{tabular}{L{0.18\textwidth}L{0.33\textwidth}L{0.39\textwidth}}
\toprule
评审维度 & 必答问题 & 最低证据 \\
\midrule
用户合同 & 谁获得或失去什么能力 & 用户旅程、权限差异、通知样例 \\
社区权力 & 版主、创作者、成员和平台权力怎样变化 & 决策矩阵、申诉和退出路径 \\
经济模型 & 谁承担成本，谁获得收入或数据价值 & 成本区间、免费层、分成与例外 \\
安全治理 & 新功能会放大哪些滥用与外部性 & 威胁模型、审核工具、升级流程 \\
基础设施 & 峰值、故障和回滚时保留哪些动作 & 容量测试、降级顺序、恢复演练 \\
数据生命周期 & 内容删除、授权撤销后衍生数据如何处理 & 保留期限、缓存清理、审计记录 \\
\bottomrule
\end{tabular}
\caption{任何高影响平台变更都应同时通过产品、治理、经济、基础设施与数据评审。}
\end{table}

\begin{table}[H]
\centering
\begin{tabular}{L{0.2\textwidth}L{0.32\textwidth}L{0.38\textwidth}}
\toprule
上线阶段 & 领先指标 & 停止或回滚条件 \\
\midrule
内部测试 & 权限正确率、误封、延迟与成本 & 无法解释的权限扩大或数据泄漏 \\
小社区试点 & 采用、投诉、版主工时、申诉成功率 & 伤害集中在弱势群体或工具不足 \\
扩大范围 & 生态迁移率、故障率、收入与内容供给 & 核心供给下降、关键客户端中断 \\
稳定运行 & 长期留存、治理一致性、成本回收 & 债务持续增长且无可信修复计划 \\
事故恢复 & 恢复时间、积压、补偿与透明度 & 数据无法核对或重复伤害未控制 \\
\bottomrule
\end{tabular}
\caption{指标必须和停止条件配对，否则监控只会记录失败而不会改变决策。}
\end{table}

\begin{warningbox}{评审时最容易漏掉的三类人}
第一类是维护公共工具但没有商业合同的开发者；第二类是承担大量版务却不出现在员工花名册中的志愿版主；第三类是被数据训练或 AI 摘要影响、却不再直接访问平台的内容作者。变更若只统计付费客户和页面活跃用户，会系统性低估这些群体的贡献与损失。
\end{warningbox}

\end{document}
